\begin{figure}[ht]
\centering
\begin{tikzpicture}[>=Stealth, font=\small, node distance=1.5cm]
  \node[draw, rounded corners, fill=blue!8, minimum width=3.3cm, minimum height=0.9cm] (mf) {模形式与自守表示};
  \node[draw, rounded corners, fill=green!10, right=3.2cm of mf, minimum width=3.3cm, minimum height=0.9cm] (ec) {椭圆曲线与代数几何};
  \node[draw, rounded corners, fill=orange!14, below=1.8cm of mf, minimum width=3.3cm, minimum height=0.9cm] (gr) {Galois 表示};
  \node[draw, rounded corners, fill=yellow!16, below=1.8cm of ec, minimum width=3.3cm, minimum height=0.9cm] (lf) {$L$-函数与特殊值};
  \node[draw, rounded corners, fill=purple!10, below=1.8cm of $(gr)!0.5!(lf)$, minimum width=3.6cm, minimum height=0.9cm] (lg) {Langlands 与 $p$-adic Langlands};
  \node[draw, rounded corners, fill=red!8, below=1.7cm of lg, minimum width=4.0cm, minimum height=0.9cm] (bsd) {BSD、Serre 猜想与更远的方向};

  \draw[<->, very thick] (mf) -- node[above] {模性} (ec);
  \draw[<->, very thick] (mf) -- node[left] {Deligne 表示} (gr);
  \draw[<->, very thick] (ec) -- node[right] {Tate 模} (lf);
  \draw[<->, very thick] (gr) -- node[below] {局部--整体} (lf);
  \draw[->, very thick] (mf) -- (lg);
  \draw[->, very thick] (ec) -- (lg);
  \draw[->, very thick] (gr) -- (lg);
  \draw[->, very thick] (lf) -- (lg);
  \draw[->, very thick] (lg) -- (bsd);
\end{tikzpicture}
\caption{后记中的大图景：模性定理不是终点，而是把模形式、Galois 表示、代数几何与 $L$-函数连入 Langlands 网络的一个成功范例。}
\label{fig:ch10-langlands-network}
\end{figure}
